\originalchapter{离散随机变量、期望与方差}
\sourcelecture{8--11}

\originalsection{从实验结果到数值分布}
一次实验的结果可能是一串正反面、一个排列或若干人的选择. 我们往往只关心其中某个数值: 正面数、某张卡的位置、到场人数. 随机变量把复杂结果转换成这样的数值, 但并不改变底层实验.
\begin{definition}
随机变量是样本空间上的实值函数 $X:\Omega\to\R$, 并要求 $\{X\le t\}$ 对每个实数 $t$ 都是事件. 若存在有限或可数集 $D\subset\R$ 使 $\Prob(X\in D)=1$, 则 $X$ 是离散随机变量. 其概率质量函数和分布函数分别为
\[
p_X(x)=\Prob(X=x),\qquad F_X(t)=\Prob(X\le t)=\sum_{x\in D,\,x\le t}p_X(x).
\]
\end{definition}
有限或可数样本空间并取所有子集为事件时, 上述事件要求自动满足. 在更一般的空间上, 它保证我们写出的概率有意义. 质量函数非负, 总和为1; 反过来, 任一满足这两点的函数都可以成为某个离散变量的质量函数. 例如取 $\Omega=D$, 令 $X(x)=x$, 以 $p_X(x)$ 给点 $x$ 分配概率即可. 同一分布可以在不同样本空间实现, 所以知道分布并不等于知道所有实验细节.

\begin{example}
从 $m$ 个不同编号中等概率随机排列, 令 $X$ 为编号1的位置. 再独立掷三个均匀的四面骰, 令 $Y$ 为点数和. 求 $X$ 的分布和 $\Prob(Y=5)$.
\end{example}
\begin{solution}
固定位置 $k$ 后, 其余编号有 $(m-1)!$ 种排列, 因此 $p_X(k)=1/m$ ($1\le k\le m$). 对第二问, 写点数为 $1+u_i$. 需要 $u_1+u_2+u_3=2$, $u_i\ge0$. 此时 $u_i\le3$ 自动成立, 非负整数解有 $\binom{4}{2}=6$ 个, 总结果数为 $4^3$, 所以概率为 $3/32$.
\end{solution}

若 $D$ 为整数集, 分布函数在相邻整数之间保持不变, 在 $k$ 处增加 $p_X(k)$. 一般情况下仍有
\[
\Prob(a<X\le b)=F_X(b)-F_X(a).
\]
区间端点不能随意互换: $\Prob(X=a)$ 可能为正. 分布函数单调不减, 右连续, 且在正负无穷处分别趋于1与0. 右连续性来自事件 $\{X\le t+1/n\}$ 递减到 $\{X\le t\}$ 和概率的连续性.

\begin{definition}
事件 $A$ 的指示变量为 $\ind_A(\omega)=1$ ($\omega\in A$), 否则为0. 因而 $\sum_{i=1}^n\ind_{A_i}$ 数出发生的事件数.
\end{definition}
指示变量不是概率: 它随实验结果取0或1, 而概率是一个固定数. 两者的联系会在期望中出现. 对排列, 定义 $A_i$ 为编号 $i$ 恰在位置 $i$, 则固定点数为 $\sum_i\ind_{A_i}$; 即使这些事件不独立, 这个逐结果恒等式仍成立.

\originalsection{多事件独立与递推}
\begin{definition}
事件 $A_1,\ldots,A_n$ 相互独立, 是指每个非空指标子集 $J$ 都满足
\[
\Prob\left(\bigcap_{j\in J}A_j\right)=\prod_{j\in J}\Prob(A_j).
\]
离散随机变量相互独立, 是指联合质量函数等于边缘质量函数的乘积.
\end{definition}
只有每一对事件满足乘法规则, 称为两两独立. 三个以上事件时, 它弱于相互独立. 例如取两个独立均匀的0、1变量 $U,V$, 事件 $U=1$, $V=1$, $U\ne V$ 都有概率 $1/2$, 每一对交集有概率 $1/4$, 但三者交集为空. 独立事件的补事件也保留相互独立性: 将一个 $A$ 换成 $A^c$ 时, 利用 $\Prob(B\cap A^c)=\Prob(B)-\Prob(B\cap A)$ 得到所需乘积, 再逐个替换. 因而独立试验中, 已知若干试验的具体结果不会改变其余试验的分布, 只要所条件化的事件概率非零.

许多概率无需枚举全部样本点, 而可按第一步分解. 第一轮后问题仍具有同样的形式时, 这种分解给出递推关系.
\begin{example}
两个参与者每轮公平地转移1单位资金, 总资金为整数 $L\ge2$. 第一人初有 $i$ 单位, $0\le i\le L$. 当一人破产时停止. 求第一人先达到 $L$ 的概率及比赛轮数的期望.
\end{example}
\begin{solution}
先确认停止不是额外假设. 从任何未停止状态出发, 连续 $L$ 轮向上就足以到达边界, 此事件概率至少 $2^{-L}$. 每一块 $L$ 轮独立于之前的试验, 因此停止时间 $\tau$ 满足
\[
\Prob(\tau>kL)\le(1-2^{-L})^k.
\]
尾概率趋于0且可求和, 所以几乎必然停止, 并且 $\E\tau<\infty$.

令 $h_i$ 为胜率. 第一轮条件分解给出 $h_i=(h_{i-1}+h_{i+1})/2$, 边界 $h_0=0,h_L=1$. 所以相邻差 $h_{i+1}-h_i$ 相等, 由边界得到 $h_i=i/L$. 破产概率为 $(L-i)/L$.

令 $t_i$ 为剩余轮数的期望. 由于期望有限, 可写
\[
t_i=1+\tfrac12t_{i-1}+\tfrac12t_{i+1},\qquad t_0=t_L=0.
\]
试取二次函数以产生常数二阶差, 代入得到 $t_i=i(L-i)$. 这也是唯一解: 两个解之差满足前面的齐次递推且边界为0, 故相邻差恒定并全为0.
\end{solution}
如果上方没有资本边界, 一个初有 $i$ 单位的公平赌徒仍以概率1最终破产. 因为在 $0$ 与任意 $L>i$ 之间比赛时, 先破产概率为 $1-i/L$; 这一事件蕴含无上界游戏最终破产. 让 $L\to\infty$ 即得概率至少为1. 这不说明平均破产时间有限; 事实上它是无穷大. 有限边界停止时间不超过无上界破产时间, 前者期望 $i(L-i)$ 随 $L$ 无界.

\begin{example}
独立试验成功概率 $p\in(0,1)$, 求在出现 $m$ 次失败之前出现 $r$ 次成功的概率, $r,m\ge1$.
\end{example}
\begin{solution}
令 $H_{r,m}$ 表示该概率, 第一轮分解得到
\[
H_{r,m}=pH_{r-1,m}+(1-p)H_{r,m-1},\quad H_{0,m}=1,\quad H_{r,0}=0.
\]
也可以预先补足到 $r+m-1$ 轮. 在这些轮次中至少有 $r$ 次成功, 当且仅当第 $r$ 次成功早于第 $m$ 次失败. 因为总轮数不足以同时包含二者, 其中一方必先达到目标. 因而
\[
H_{r,m}=\sum_{k=r}^{r+m-1}\binom{r+m-1}{k}p^k(1-p)^{r+m-1-k}.
\]
公平情况下这就是 Pascal 三角形的一段尾和除以 $2^{r+m-1}$. 它也满足上面的递推, 将计数公式与逐步分析联系起来.
\end{solution}

\originalsection{期望、函数与收敛条件}
\begin{definition}
离散随机变量 $X$ 的有限期望定义为
\[
\E X=\sum_x xp_X(x),\qquad\text{条件是 }\sum_x|x|p_X(x)<\infty.
\]
此时称 $X$ 可积. 对非负变量也允许期望为 $+\infty$, 其值是有限部分和的上确界.
\end{definition}
绝对收敛条件保证枚举顺序不影响答案. 有正负项时, 不可用一个条件收敛级数定义概率论中的有限期望. 若正部分和负部分的期望都为无穷, 表达式 $+\infty-\infty$ 没有定义.

期望是按概率加权的平均值, 不必是可能出现的结果. 均匀骰子的期望为 $7/2$, 尽管没有这样的点数. 在可数样本空间中, 既可按样本点求和, 也可先把产生相同 $X$ 值的样本点合并:
\[
\sum_{\omega\in\Omega}X(\omega)\Prob(\{\omega\})
=\sum_xx\sum_{\omega:X(\omega)=x}\Prob(\{\omega\})
=\sum_xxp_X(x).
\]
对非负项, 合并由有限部分和的上确界解释; 对有符号项, 绝对收敛允许重排.
\begin{proposition}[函数的期望]
若 $\E|g(X)|<\infty$, 则 $\E g(X)=\sum_xg(x)p_X(x)$. 若 $g\ge0$, 公式在允许无穷的意义下仍成立.
\end{proposition}
\begin{proof}
设 $g(X)$ 取值 $y$. 它的质量函数为 $\sum_{x:g(x)=y}p_X(x)$. 将 $\sum_yy\Prob(g(X)=y)$ 依这一分组展开, 得到所述公式; 重排的依据如上.
\end{proof}
这避免了每次先求 $g(X)$ 的分布. 例如骰子平方的期望为 $\sum_{j=1}^6j^2/6=91/6$, 并不等于 $(7/2)^2$.

\begin{theorem}[期望的线性性]
若 $X_1,\ldots,X_n$ 可积, $a_1,\ldots,a_n,b$ 为实常数, 则
\[
\E\left(b+\sum_{i=1}^na_iX_i\right)=b+\sum_{i=1}^na_i\E X_i.
\]
不要求随机变量独立.
\end{theorem}
\begin{proof}
二变量情形取联合质量函数 $p(x,y)$, 有
\[
\sum_{x,y}|ax+by|p(x,y)\le |a|\E|X|+|b|\E|Y|<\infty.
\]
因此可以分拆和式并分别按另一个变量求和:
$\sum_{x,y}(ax+by)p(x,y)=a\sum_xxp_X(x)+b\sum_yyp_Y(y)$.
常数期望等于常数, 再有限次应用即得结论.
\end{proof}
特别地, $\E\ind_A=\Prob(A)$. 求一个计数变量的期望时, 往往无需知道整个分布.
\begin{example}
将 $n$ 个编号均匀随机排列, 固定点数为 $K$. 求 $\E K$.
\end{example}
\begin{solution}
$K=\sum_{i=1}^n\ind_{A_i}$, 每个编号保持原位的概率为 $(n-1)!/n!=1/n$. 所以 $\E K=\sum_i1/n=1$. 这不是说每次恰有一个固定点; 某些排列没有, 有些有很多.
\end{solution}

\begin{proposition}[整数变量的尾和公式]
若 $X$ 取非负整数值, 则
\[
\E X=\sum_{k=1}^{\infty}\Prob(X\ge k),
\]
两边可同时为无穷.
\end{proposition}
\begin{proof}
逐结果有 $X=\sum_{k\ge1}\ind_{\{X\ge k\}}$. 等价地, 直接展开
$\sum_{j\ge0}jp_X(j)=\sum_{j\ge0}\sum_{k=1}^jp_X(j)$.
所有项非负, 双重和等于各有限子和的上确界, 因而可交换求和顺序, 得到 $\sum_{k\ge1}\sum_{j\ge k}p_X(j)$.
\end{proof}
有限个变量的线性性不能无条件用于有符号无穷和. 非负项可按上述理由使用; 有符号项则需例如 $\sum_i\E|X_i|<\infty$ 来保证绝对可积.

\originalsection{均值之外的波动}
相同的均值可能对应很不同的结果. 恒付1与以概率 $1/M$ 付 $M$、否则付0都有均值1, 后者却随 $M$ 增大而越来越不稳定. 期望描述中心, 方差描述偏离中心的平方大小.
\begin{definition}
若 $\E X^2<\infty$, 定义
\[
\Var(X)=\E[(X-\E X)^2],\qquad
\operatorname{SD}(X)=\sqrt{\Var(X)}.
\]
\end{definition}
二阶矩有限保证一阶矩有限, 因为 $|x|\le1+x^2$. 平方消除了正负偏离相互抵消的可能, 也加重了极端值的影响. 方差的量纲是原量纲的平方, 标准差则与变量相同.
\begin{proposition}
二阶矩有限时,
\[
\Var(X)=\E X^2-(\E X)^2,\quad
\Var(aX+b)=a^2\Var(X),\quad
\operatorname{SD}(aX+b)=|a|\operatorname{SD}(X).
\]
且 $\Var(X)=0$ 当且仅当 $X=\E X$ 几乎必然成立.
\end{proposition}
\begin{proof}
设 $\mu=\E X$, 展开平方并使用线性性得
$\E(X-\mu)^2=\E X^2-2\mu\E X+\mu^2=\E X^2-\mu^2$.
又 $aX+b-\E(aX+b)=a(X-\mu)$, 得到缩放公式. 若非负变量 $(X-\mu)^2$ 的期望为0, 则每个正概率取值都只能为0, 故 $X=\mu$ 几乎必然; 反向直接成立.
\end{proof}
取绝对值不可省去: $a=-1$ 不会使标准差变成负数. 对均匀骰子, 方差为 $91/6-49/4=35/12$. 对上述稀有支付, 二阶矩为 $M$, 方差为 $M-1$. 这解释了为什么只有均值不能判断风险.

\begin{definition}
二阶矩有限的 $X,Y$ 的协方差为
\[
\Cov(X,Y)=\E[(X-\E X)(Y-\E Y)]=\E(XY)-\E X\E Y.
\]
\end{definition}
积可积, 因为 $2|XY|\le X^2+Y^2$. 协方差的符号描述中心化变量是否倾向于同向变化.
\begin{proposition}[和的方差]
\[
\Var\left(\sum_i a_iX_i\right)=\sum_i a_i^2\Var(X_i)
+2\sum_{i<j}a_ia_j\Cov(X_i,X_j).
\]
独立变量的协方差为0, 因而独立变量的方差可以相加.
\end{proposition}
\begin{proof}
将 $\sum_i a_i(X_i-\E X_i)$ 的平方展开并求期望即可. 独立时, 联合质量函数分解, 而绝对可积允许求和, 得 $\E XY=(\E X)(\E Y)$.
\end{proof}
零协方差不保证独立. 令 $X$ 均匀取 $-1,0,1$, $Y=X^2$. 则 $\E X=\E X^3=0$, 故协方差为0; 但知道 $Y=0$ 就确定 $X=0$.

\begin{example}
对 $n\ge2$ 个编号的随机排列, 求固定点数 $K$ 的方差.
\end{example}
\begin{solution}
设 $I_i=\ind_{A_i}$. 平方展开给出
\[
\E K^2=\sum_i\E I_i^2+\sum_{i\ne j}\E I_iI_j.
\]
$I_i^2=I_i$, 贡献为 $n/n=1$. 同时固定两个指定编号有 $(n-2)!$ 种排列, 概率为 $1/[n(n-1)]$, $n(n-1)$ 个有序非对角项再贡献1. 所以 $\E K^2=2$, $\Var(K)=1$. 若 $n=1$, $K\equiv1$, 方差为0, 不能使用含 $n-1$ 分母的公式.
\end{solution}

\originalsection{Bernoulli试验与二项分布}
\begin{definition}
取值0、1且 $\Prob(I=1)=p$ 的变量称为 Bernoulli 变量, $0\le p\le1$. 独立的 $n$ 个同参数 Bernoulli 变量之和服从二项分布, 记作 $X\sim\Bin(n,p)$, $n$ 为非负整数.
\end{definition}
令 $q=1-p$. 恰有 $k$ 次成功需指定 $k$ 个位置, 每一种序列的概率为 $p^kq^{n-k}$, 故
\[
\Prob(X=k)=\binom nkp^kq^{n-k},\qquad0\le k\le n.
\]
其他整数概率为0. 总和为 $(p+q)^n=1$. 在 $p=0$ 或1时直接理解为退化分布, 避免依赖 $0^0$ 记号. Bernoulli 变量有 $\E I=p$, $\Var(I)=p-p^2=pq$.

\begin{theorem}
$X\sim\Bin(n,p)$ 时, $\E X=np$, $\Var(X)=npq$. 若 $n\ge1$, $Y\sim\Bin(n-1,p)$, 则对整数 $r\ge1$,
\[
\E X^r=np\E(Y+1)^{r-1}.
\]
\end{theorem}
\begin{proof}
指示变量分解和线性性给出均值 $np$, 独立性给出方差 $npq$. 为说明直接求和如何工作, 使用
$k\binom nk=n\binom{n-1}{k-1}$:
\begin{align*}
\E X^r
&=\sum_{k=1}^nk^{r-1}k\binom nkp^kq^{n-k}\\
&=np\sum_{j=0}^{n-1}(j+1)^{r-1}\binom{n-1}{j}p^jq^{n-1-j}.
\end{align*}
这正是所述矩递推. $r=1$ 时剩余和为1; $r=2$ 时 $\E X^2=np[(n-1)p+1]$, 减去 $(np)^2$ 再得方差. 边界参数可直接检查.
\end{proof}
在均值计算中, 每项多出的 $k$ 不是障碍, 而是使 Pascal 三角形的第 $n$ 行转成第 $n-1$ 行的线索. 更一般地, 下降阶乘 $(X)_r=X(X-1)\cdots(X-r+1)$ 满足
\[
\E(X)_r=(n)_rp^r\quad(1\le r\le n).
\]
因为 $(X)_r$ 数出成功试验中的有序 $r$ 元组, 一共有 $(n)_r$ 个候选, 各个同时成功概率为 $p^r$. 这也可以用于高阶矩计算.

\begin{example}
一场活动有24个座位, 发出26份确认, 每人独立以概率 $9/10$ 到场. 求座位不足的概率. 另有12人独立公平选择两项方案, 求两方人数相同的概率.
\end{example}
\begin{solution}
到场人数 $A\sim\Bin(26,9/10)$, 座位不足是 $A\ge25$, 因此概率为
\[
\binom{26}{25}(9/10)^{25}(1/10)+(9/10)^{26}.
\]
也可数缺席者 $B=26-A\sim\Bin(26,1/10)$, 条件是 $B\le1$, 得到同式. 第二问为 $\binom{12}{6}2^{-12}$. 固定赞成者或固定反对者若另外存在, 应先扣除其人数, 再确定随机部分需要的成功数.
\end{solution}
应用前必须核对固定试验次数、相同成功概率和独立性. 不放回抽样通常不满足独立性, 下一章将处理; 不同人的到场概率不同也不再是普通二项分布, 尽管期望仍为各概率之和.

\originalsection{期望效用与投资组合}
期望的用途之一是长期平均, 后面的极限定理会说明这一联系的精确条件. 另一个用途是比较行动结果. 但最大化金钱期望与最大化个人满意度并非同一问题.
\begin{example}
某人有初始财富10. 方案甲确定再得2, 方案乙以各 $1/2$ 的概率再得0或4. 以 $u(w)=\sqrt w$ 表示财富效用, 比较两方案.
\end{example}
\begin{solution}
两方案最终财富期望都为12. 甲的期望效用为 $\sqrt{12}$, 乙为 $(\sqrt{10}+\sqrt{14})/2$. 平方比较或利用平方根的严格凹性可得后者较小. 直接计算:
\[
\left(\frac{\sqrt{10}+\sqrt{14}}2\right)^2=6+\frac{\sqrt{140}}2<12,
\]
因为 $140<144$. 因而按此效用偏好确定支付. 效用函数和财富定义必须先指定; 不能仅凭均值断言任何人都应作同样选择.
\end{solution}
这里使用通常约定: 效用越大越好. 若使用损失或成本函数, 应最小化其期望. 这两种方向不能混用.

投资组合中, 收益的线性性与波动的交叉项尤其清楚. 设两个资产的单位收益为 $R_1,R_2$, 权重 $w$ 与 $1-w$, 则组合收益为 $R=wR_1+(1-w)R_2$. 在二阶矩有限时,
\[
\E R=w\mu_1+(1-w)\mu_2,
\quad\Var(R)=w^2v_1+(1-w)^2v_2+2w(1-w)c,
\]
其中 $\mu_i=\E R_i$, $v_i=\Var(R_i)$, $c=\Cov(R_1,R_2)$. 对 $0\le w\le1$ 是不卖空的模型; 若允许实数权重, 则是在相应理想化借贷与卖空约束下计算.
\begin{example}
两个资产收益均值相同, 方差都为4. 比较各投一半时协方差为0、4、$-4$ 的组合方差.
\end{example}
\begin{solution}
等权组合方差为 $\tfrac14\cdot4+\tfrac14\cdot4+\tfrac12c=2+c/2$. 三种情形分别为2、4、0. 它们可以实现: 独立同分布的中心化收益实现 $c=0$, 相同收益实现 $c=4$, 相反中心化收益实现 $c=-4$. 第三种情形两部分波动逐结果抵消. 因此资产数量本身不能说明分散程度, 相关性同样重要.
\end{solution}
若 $v_1+v_2-2c=\Var(R_1-R_2)>0$, 将方差作为 $w$ 的二次函数求导得到无约束最小点
\[
w_* =\frac{v_2-c}{v_1+v_2-2c}.
\]
若要求 $0\le w\le1$, 应将此点截到区间内. 分母为0时 $R_1-R_2$ 几乎必然为常数, 各权重的方差相同, 不可除以0. 金融合约的无套利定价有时使用另一种风险中性概率下的期望; 它不是实际发生概率下的期望, 其条件和构造将在金融模型章节给出.

\originalsection{练习与完整解答}
\begin{exercise}
随机变量 $X$ 以概率 $1/4,1/2,1/4$ 分别取 $-2,1,4$. 求 $F_X(t)$、$\E X$、$\Var(X)$ 及 $\E(2X-1)^2$.
\end{exercise}
\begin{solution}
分布函数在 $t<-2$ 时为0, 在 $-2\le t<1$ 时为 $1/4$, 在 $1\le t<4$ 时为 $3/4$, 在 $t\ge4$ 时为1. $\E X=-1/2+1/2+1=1$, $\E X^2=1+1/2+4=11/2$, 故方差为 $9/2$. 最后一项为 $4\E X^2-4\E X+1=19$.
\end{solution}
\begin{exercise}
独立的 $n$ 次试验成功概率为 $p$. 令 $A$ 是成功数, $B$ 是失败数. 求二者的协方差, 并解释为何它们一般不独立.
\end{exercise}
\begin{solution}
$B=n-A$, 所以 $\Cov(A,B)=\Cov(A,n-A)=-\Var(A)=-np(1-p)$. 当 $n\ge1$ 且 $0<p<1$, 它为负, 因而不独立. 更直接地, $A$ 的数值完全确定 $B$. 在退化参数下两者都是常量, 可以独立.
\end{solution}
\begin{exercise}
均匀随机排列中, 令 $C$ 为位置1、2是否都保持原位的指示变量 ($n\ge2$). 求 $\E C$, 并比较 $\E(\ind_{A_1}\ind_{A_2})$ 与 $\E\ind_{A_1}\E\ind_{A_2}$.
\end{exercise}
\begin{solution}
$\E C=1/[n(n-1)]$, 而边缘期望乘积为 $1/n^2$. 前者较大, 表明一个编号固定后, 另一个固定的条件概率从 $1/n$ 增为 $1/(n-1)$. 此处不可以套用独立 Bernoulli 和的方差.
\end{solution}
